% GENERATED FILE. Edit canonical lessons, then run npm run generate:books.
% canonical-source-hash: fnv1a64:67a362f1881a4b18
% canonical-lessons: JA-C356-jouken, JA-C356-tadashi, JA-C356-angai, JA-C356-taitei, JA-C356-itsumademo, JA-R356-review-352-356

\chapter{On One Condition}
\label{ch:ja-on-one-condition}
\hlchaptermodality{356}

\begin{chapteropening}
\textbf{By the end of this chapter:} \emph{I can set a condition and say what happens if it is met.}

Complete the last lesson of chapter 356: I can set a condition and say what happens if it is met.
\end{chapteropening}

\section[\texorpdfstring{\ja{じょうけん} (jōken)}{じょうけん (jōken)}]{\ja{じょうけん} (jōken) — a condition; terms}

\label{lesson:JA-C356-jouken}



\begin{quote}
Before the new one: say the Japanese for a possibility, then the Japanese for to miss the mark.
\end{quote}

\subsection*{You'll want to know: \ja{じょうけん}}
\textbf{\ja{じょうけん}} — \emph{jōken} — \textquotedblleft{}a condition; terms\textquotedblright{}.

\begin{quote}
\textbf{\ja{じょうけんが} \ja{あったら、}\ja{言}\ja{って} \ja{ください。}} — \emph{jōken ga attara, itte kudasai} — \textquotedblleft{}If you have any conditions, please say so.\textquotedblright{}
\end{quote}

\begin{cousinweb}
\textbf{\ja{じょうけん}} (\emph{jōken}) is Sino-Japanese (kango): two Chinese characters read with their on'yomi, readings borrowed from Chinese. They are \emph{jō}, \textquotedblleft{}article\textquotedblright{} and \emph{ken}, \textquotedblleft{}matter\textquotedblright{} (Middle Chinese \emph{gjenX}) (KANJIDIC2; Middle Chinese in Baxter's transcription).
\end{cousinweb}

\begin{grammarlens}[title={What you've built}]
One more word for saying what happens if.
\end{grammarlens}

\subsection*{Your turn}
\begin{itemize}
\raggedright
  \item \emph{Say it:} \emph{jōken}
  \item \emph{Say it:} \emph{jōken}, once more
  \item \emph{Say it:} \emph{jōken ga attara, itte kudasai} — If you have any conditions, please say so.
  \item \emph{Recall:} say the Japanese for a possibility, then the Japanese for to miss the mark, then say \emph{jōken} again
\end{itemize}

\begin{tcolorbox}[breakable,colback=teal!4,colframe=teal!35!black,title={Before you move on}]
What does \ja{じょうけん} mean? (\textquotedblleft{}A condition; terms\textquotedblright{}.) Say it once more.
\end{tcolorbox}
\section[\texorpdfstring{\ja{ただし} (tadashi)}{ただし (tadashi)}]{\ja{ただし} (tadashi) — but; however; provided that}

\label{lesson:JA-C356-tadashi}



\begin{quote}
Before the new one: say the Japanese for to miss the mark, then the Japanese for a condition.
\end{quote}

\subsection*{You'll want to know: \ja{ただし}}
\textbf{\ja{ただし}} — \emph{tadashi} — \textquotedblleft{}but; however; provided that\textquotedblright{}.

\begin{quote}
\textbf{\ja{ただし、ねつが} \ja{あったら、やすんで} \ja{ください。}} — \emph{tadashi, netsu ga attara, yasunde kudasai} — \textquotedblleft{}But if you have a fever, take a rest.\textquotedblright{}
\end{quote}

\begin{cousinweb}
\textbf{\ja{ただし}} (\emph{tadashi}) is a native Japanese word (wago), inherited rather than borrowed from Chinese, and its older origin is not settled. JMdict notes that it is usually written in kana, and KANJIDIC2 gives \emph{tadashi} as a native reading of its kanji (JMdict, KANJIDIC2).
\end{cousinweb}

\begin{grammarlens}[title={What you've built}]
One more word for saying what happens if.
\end{grammarlens}

\subsection*{Your turn}
\begin{itemize}
\raggedright
  \item \emph{Say it:} \emph{tadashi}
  \item \emph{Say it:} \emph{tadashi}, once more
  \item \emph{Say it:} \emph{tadashi, netsu ga attara, yasunde kudasai} — But if you have a fever, take a rest.
  \item \emph{Recall:} say the Japanese for to miss the mark, then the Japanese for a condition, then say \emph{tadashi} again
\end{itemize}

\begin{tcolorbox}[breakable,colback=teal!4,colframe=teal!35!black,title={Before you move on}]
What does \ja{ただし} mean? (\textquotedblleft{}But; however; provided that\textquotedblright{}.) Say it once more.
\end{tcolorbox}
\section[\texorpdfstring{\ja{あんがい} (angai)}{あんがい (angai)}]{\ja{あんがい} (angai) — unexpectedly; surprisingly}

\label{lesson:JA-C356-angai}



\begin{quote}
Before the new one: say the Japanese for a condition, then the Japanese for provided that.
\end{quote}

\subsection*{You'll want to know: \ja{あんがい}}
\textbf{\ja{あんがい}} — \emph{angai} — \textquotedblleft{}unexpectedly; surprisingly\textquotedblright{}.

\begin{quote}
\textbf{\ja{はじめたら、あんがい} \ja{かんたんでしょう。}} — \emph{hajimetara, angai kantan deshō} — \textquotedblleft{}Once you start, it will probably be surprisingly easy.\textquotedblright{}
\end{quote}

\begin{cousinweb}
\textbf{\ja{あんがい}} (\emph{angai}) is Sino-Japanese (kango): two Chinese characters read with their on'yomi, readings borrowed from Chinese. They are \emph{an}, \textquotedblleft{}plan\textquotedblright{} (Middle Chinese \emph{'anH}) and \emph{gai}, \textquotedblleft{}outside\textquotedblright{} (\emph{ngwajH}) (KANJIDIC2; Middle Chinese in Baxter's transcription).
\end{cousinweb}

\begin{grammarlens}[title={What you've built}]
One more word for saying what happens if.
\end{grammarlens}

\subsection*{Your turn}
\begin{itemize}
\raggedright
  \item \emph{Say it:} \emph{angai}
  \item \emph{Say it:} \emph{angai}, once more
  \item \emph{Say it:} \emph{hajimetara, angai kantan deshō} — Once you start, it will probably be surprisingly easy.
  \item \emph{Recall:} say the Japanese for a condition, then the Japanese for provided that, then say \emph{angai} again
\end{itemize}

\begin{tcolorbox}[breakable,colback=teal!4,colframe=teal!35!black,title={Before you move on}]
What does \ja{あんがい} mean? (\textquotedblleft{}Unexpectedly; surprisingly\textquotedblright{}.) Say it once more.
\end{tcolorbox}
\section[\texorpdfstring{\ja{たいてい} (taitei)}{たいてい (taitei)}]{\ja{たいてい} (taitei) — in most cases; usually}

\label{lesson:JA-C356-taitei}



\begin{quote}
Before the new one: say the Japanese for provided that, then the Japanese for unexpectedly.
\end{quote}

\subsection*{You'll want to know: \ja{たいてい}}
\textbf{\ja{たいてい}} — \emph{taitei} — \textquotedblleft{}in most cases; usually\textquotedblright{}.

\begin{quote}
\textbf{\ja{あめが} \ja{ふったら、たいてい} \ja{うちに} \ja{います。}} — \emph{ame ga futtara, taitei uchi ni imasu} — \textquotedblleft{}When it rains, I usually stay at home.\textquotedblright{}
\end{quote}

\begin{cousinweb}
\textbf{\ja{たいてい}} (\emph{taitei}) is Sino-Japanese (kango): two Chinese characters read with their on'yomi, readings borrowed from Chinese. They are \emph{tai}, \textquotedblleft{}large\textquotedblright{} (Middle Chinese \emph{dajH}) and \emph{tei}, \textquotedblleft{}reach\textquotedblright{} (\emph{tejX}) (KANJIDIC2; Middle Chinese in Baxter's transcription). JMdict notes that it is usually written in kana.
\end{cousinweb}

\begin{grammarlens}[title={What you've built}]
One more word for saying what happens if.
\end{grammarlens}

\subsection*{Your turn}
\begin{itemize}
\raggedright
  \item \emph{Say it:} \emph{taitei}
  \item \emph{Say it:} \emph{taitei}, once more
  \item \emph{Say it:} \emph{ame ga futtara, taitei uchi ni imasu} — When it rains, I usually stay at home.
  \item \emph{Recall:} say the Japanese for provided that, then the Japanese for unexpectedly, then say \emph{taitei} again
\end{itemize}

\begin{tcolorbox}[breakable,colback=teal!4,colframe=teal!35!black,title={Before you move on}]
What does \ja{たいてい} mean? (\textquotedblleft{}In most cases; usually\textquotedblright{}.) Say it once more.
\end{tcolorbox}
\section[\texorpdfstring{\ja{いつまでも} (itsumademo)}{いつまでも (itsumademo)}]{\ja{いつまでも} (itsumademo) — forever; for as long as you like}

\label{lesson:JA-C356-itsumademo}



\begin{quote}
Before the new one: say the Japanese for unexpectedly, then the Japanese for in most cases.
\end{quote}

\subsection*{You'll want to know: \ja{いつまでも}}
\textbf{\ja{いつまでも}} — \emph{itsumademo} — \textquotedblleft{}forever; for as long as you like\textquotedblright{}.

\begin{quote}
\textbf{\ja{じかんが} \ja{あったら、いつまでも} \ja{いて} \ja{ください。}} — \emph{jikan ga attara, itsumademo ite kudasai} — \textquotedblleft{}If you have time, stay as long as you like.\textquotedblright{}
\end{quote}

\begin{cousinweb}
\textbf{\ja{いつまでも}} (\emph{itsumademo}) is a native Japanese word (wago), inherited rather than borrowed from Chinese, and its older origin is not settled. JMdict notes that it is usually written in kana (JMdict).
\end{cousinweb}

\begin{grammarlens}[title={What you've built}]
One more word for saying what happens if.
\end{grammarlens}

\subsection*{Your turn}
\begin{itemize}
\raggedright
  \item \emph{Say it:} \emph{itsumademo}
  \item \emph{Say it:} \emph{itsumademo}, once more
  \item \emph{Say it:} \emph{jikan ga attara, itsumademo ite kudasai} — If you have time, stay as long as you like.
  \item \emph{Recall:} say the Japanese for unexpectedly, then the Japanese for in most cases, then say \emph{itsumademo} again
\end{itemize}

\begin{tcolorbox}[breakable,colback=teal!4,colframe=teal!35!black,title={Before you move on}]
What does \ja{いつまでも} mean? (\textquotedblleft{}Forever; for as long as you like\textquotedblright{}.) Say it once more.
\end{tcolorbox}
\section[(dialogue)]{Looking back: from Little Accidents to On One Condition}

\label{lesson:JA-R356-review-352-356}



\begin{quote}
Say the Japanese for to be spilled and to come undone.
\end{quote}

\subsection*{Your turn}
Say these aloud:

\begin{itemize}
\raggedright
  \item to be spilled, to come undone, to tip over, to slip out of place, a tool
  \item to inspect, how things are running, a fire extinguisher, a water outage, tentatively
  \item a downpour, a cold snap, dampness, a parasol, in advance
  \item perhaps, surely, very likely, a possibility, to miss the mark
  \item a condition, provided that, unexpectedly, in most cases, forever
\end{itemize}

\begin{tcolorbox}[breakable,colback=teal!4,colframe=teal!35!black,title={Before you move on}]
To be spilled? (\textbf{\ja{こぼれる}}, \emph{koboreru}.) To come undone? (\textbf{\ja{ほどける}}, \emph{hodokeru}.) To tip over? (\textbf{\ja{ひっくりかえる}}, \emph{hikkurikaeru}.) To slip out of place? (\textbf{\ja{ずれる}}, \emph{zureru}.) A tool? (\textbf{\ja{どうぐ}}, \emph{dōgu}.) To inspect? (\textbf{\ja{てんけんする}}, \emph{tenken suru}.) How things are running? (\textbf{\ja{ちょうし}}, \emph{chōshi}.) A fire extinguisher? (\textbf{\ja{しょうかき}}, \emph{shōkaki}.) A water outage? (\textbf{\ja{だんすい}}, \emph{dansui}.) Tentatively? (\textbf{\ja{いちおう}}, \emph{ichiō}.) A downpour? (\textbf{\ja{どしゃぶり}}, \emph{doshaburi}.) A cold snap? (\textbf{\ja{かんぱ}}, \emph{kanpa}.) Dampness? (\textbf{\ja{しっけ}}, \emph{shikke}.) A parasol? (\textbf{\ja{ひがさ}}, \emph{higasa}.) In advance? (\textbf{\ja{あらかじめ}}, \emph{arakajime}.) Perhaps? (\textbf{\ja{たぶん}}, \emph{tabun}.) Surely? (\textbf{\ja{きっと}}, \emph{kitto}.) Very likely? (\textbf{\ja{おそらく}}, \emph{osoraku}.) A possibility? (\textbf{\ja{かのうせい}}, \emph{kanōsei}.) To miss the mark? (\textbf{\ja{はずれる}}, \emph{hazureru}.) A condition? (\textbf{\ja{じょうけん}}, \emph{jōken}.) Provided that? (\textbf{\ja{ただし}}, \emph{tadashi}.) Unexpectedly? (\textbf{\ja{あんがい}}, \emph{angai}.) In most cases? (\textbf{\ja{たいてい}}, \emph{taitei}.) Forever? (\textbf{\ja{いつまでも}}, \emph{itsumademo}.)
\end{tcolorbox}
